%%%%%%%%%%  scan idx 134  |  folio 121  %%%%%%%%%%
\SGAPrintPage{121}%
when \(F\) is strictly perfect, the arrow \(\Id_F\) is locally cofinal in \((\underline{\Qis/F})\).

\sgaStableTarget{sga6.e01.notation.4-5}\textbf{Notation 4.5.} Let \(\underline{M}\) be a fibered \(\underline{S}\)-category. Let \(X\in\Ob \underline{S}\), and \(E,F\in\Ob\underline{M}_X\). We shall denote by \(\underline{\Hom}_{\underline{M}}(E,F)\) the sheaf on \(\underline{S}_{/X}\) associated to the presheaf of \(\underline{S}\)-morphisms from \(E\) to \(F\) in \(\underline{M}\) (for the definition of this presheaf, see [G] I 2.6; take care that, even when \(\underline{C}\) is a stack ([G]), the corresponding presheaf for \(\underline{M}=\Kpx(\underline{C})\) (or \(\Dpx(\underline{C})\)) is in general not a sheaf).

\sgaStableTarget{sga6.e01.cor.4-6}\textbf{Corollary 4.6.} \uline{Let} \(E,F\in\Ob\Cpx(\underline{C}_S)\)\uline{, with} \(E\) \uline{strictly perfect. The canonical arrow}
\[
  \underline{\Hom}_{\Kpx(\underline{C})}(E,F)
  \longrightarrow
  \underline{\Hom}_{\Dpx(\underline{C})}(E,F)
\]
\uline{is an isomorphism}.

\uline{Proof.} This is an immediate consequence of (4.4) and of the formula
\[
  \Hom_{\Dpx(\underline{C}_X)}(E_X,F_X)
  = \varinjlim_{(\Qis/E_X)^\circ}\Hom_{\Kpx(\underline{C}_X)}(\,.\,,F_X)
\]
for \(X\) over \(S\).

\sgaStableTarget{sga6.e01.def.4-7}\textbf{Definition 4.7.} \uline{Let} \(F\in\Ob\Dpx(\underline{C}_S)\) \uline{and let} \(a,b\in\mathbb Z\) \uline{be such that} \(a\leq b\)\uline{. We shall say that} \(F\) \uline{has} \(\underline{C}_0\)-perfect amplitude contained in the interval \([a,b]\), \uline{and we shall write}
\[
  \underline{C}_0\text{-}\mathrm{parf}\text{-}\amp(F)\subset[a,b],
\]
\uline{when the set of} \(X\) \uline{over} \(S\) \uline{for which there exists a quasi-isomorphism}
%%%%%%%%%%  scan idx 135  |  folio 122  %%%%%%%%%%
\SGAPrintPage{122}%
\(F'\to F\) \uline{of} \(\Cpx(\underline{C}_X)\) \uline{with} \(F'\) \uline{strictly} \(\underline{C}_0\)\uline{-perfect} (2.1)\uline{, and such that} \(F'^i=0\) \uline{for} \(i\notin[a,b]\)\uline{, is a refinement of} \(S\). (\uline{By} (4.3) \uline{this condition is equivalent to the one obtained by replacing the expression ``quasi-isomorphism of} \(\Cpx(\underline{C}_X)\)\uline{'' by ``isomorphism of} \(\Dpx(\underline{C}_X)\)\uline{''.}) \uline{We shall say that} \(F\) \uline{has} finite \(\underline{C}_0\)-perfect amplitude \uline{when there exist integers} \(a\) \uline{and} \(b\) \uline{such that} \(\underline{C}_0\)-\(\mathrm{parf}\)-\(\amp(F)\subset[a,b]\). \uline{We shall say that} \(F\) \uline{is} \(\underline{C}_0\)-perfect \uline{when} \(F\) \uline{locally has finite} \(\underline{C}_0\)\uline{-perfect amplitude}\SGAAuthorFootnote{raisedstar}{Of course, when no confusion can arise, we shall omit \(\underline{C}_0\) from the preceding expressions.}.

\sgaStableTarget{sga6.e01.ex.4-8}\textbf{Example 4.8.} Let \((\underline{S},A)\) be a ringed topos, \(\underline{C}\) the \(\underline{S}\)-category defined by \(U\longmapsto\) (category of left \(A_U\)-Modules), and \(\underline{C}_0\) the sub-\(\underline{S}\)-category defined by \(U\longmapsto\) (category of left \(A_U\)-Modules that are direct factors of free Modules of finite type). By (1.3.1), the pair \((\underline{C},\underline{C}_0)\) satisfies the hypotheses (4.0). An object \(F\) of \(\Dpx(_A \underline{S})\) (cf.\ (2.15)) will be said to have perfect amplitude contained in \([a,b]\) (resp.\ finite perfect amplitude, resp.\ to be perfect) when \(F\) has \(\underline{C}_0\)-perfect amplitude contained in \([a,b]\) (resp.\dots). When the direct factors of free Modules of finite type are locally free (cf.\ (2.15.1)), saying that \(F\) is perfect therefore means that \(F\) is locally isomorphic (in \(\Dpx(\underline{S})\)) to a complex with bounded degrees and components free of finite type.

\sgaStableTarget{sga6.e01.notation.4-9}\textbf{Notation 4.9.} We shall denote by \(\Dpx(\underline{C})_{\underline{C}_0\text{-}\mathrm{parf}}\) or \(\Parf(\underline{C},\underline{C}_0)\) the strictly full sub-\(\underline{S}\)-category of \(\Dpx(\underline{C})\) formed by the \(\underline{C}_0\)-perfect objects. In the case of Example (4.8), we shall write \(\Dpx(_A \underline{S})_{\mathrm{parf}}\), or \(\Parf(_A \underline{S})\), or simply \(\Parf(\underline{S})\) when there is no possible ambiguity about \(A\).

%%%%%%%%%%  scan idx 136  |  folio 123  %%%%%%%%%%
\SGAPrintPage{123}%
\sgaStableTarget{sga6.e01.prop.4-10}\textbf{Proposition 4.10.} \uline{The category} \(\underline{\Parf}(\underline{C},\underline{C}_0)\) \uline{is a triangulated sub-}\(\underline{S}\)\uline{-category of} \(\Dpx(\underline{C})\) (1.1.4)\uline{. The canonical functor}
\[
  \Kpx^b(\underline{C}_0)\longrightarrow\Parf(\underline{C},\underline{C}_0)
\]
(resp.\ \(\Tr(\Kpx^b(\underline{C}_0))\to\Tr(\Parf(\underline{C},\underline{C}_0))\) (1.1.3)) \uline{induces an} \(\underline{S}\)\uline{-equivalence on the associated stacks} ([G] II 2.2.4)\uline{, which means}:
\begin{enumerate}[label=(\roman*)]
\item \uline{for every} \(X\in\Ob \underline{S}\) \uline{and all} \(E,F\in\Ob\Kpx^b(\underline{C}_{0X})\) (resp.\ \(\Tr(\Kpx^b(\underline{C}_{0X}))\)), \uline{the canonical arrow}
\[
  \underline{\Hom}_{\Kpx^b(\underline{C}_0)}(E,F)
  \longrightarrow
  \underline{\Hom}_{\Dpx(\underline{C})}(E,F)
\]
(\uline{resp.}\ \(\underline{\Hom}_{\Tr(\Kpx^b(\underline{C}_0))}(E,F)\to\underline{\Hom}_{\Tr(\Dpx(\underline{C}))}(E,F)\)) \uline{is an isomorphism, and}
\item \uline{every object of} \(\underline{\Parf}(\underline{C},\underline{C}_0)\) (resp.\ \(\Tr(\Parf(\underline{C},\underline{C}_0))\)) \uline{is locally isomorphic in} \(\Dpx(\underline{C})\) \uline{(resp.}\ \(\Tr(\Dpx(\underline{C}))\)) \uline{to an object of} \(\Kpx^b(\underline{C}_0)\) \uline{(resp.}\ \(\Tr(\Kpx^b(\underline{C}_0))\)).
\end{enumerate}

\uline{Proof.} Assertion (i) without the ``resp.'' alternative is a special case of (4.6), and assertion (ii) without the ``resp.'' alternative is true by definition ((4.7) and (4.9)). It is also clear that \(\Parf(\underline{C},\underline{C}_0)\) is stable under translation of degrees. To prove that \(\Parf(\underline{C},\underline{C}_0)\) is a triangulated sub-\(\underline{S}\)-category of \(\Dpx(\underline{C})\), it therefore suffices to show that, when
\sgaStableTarget{sga6.e01.eq.unnamed.occ3}\[
\begin{tikzcd}
& G \arrow[dl,"+1"'] & \\
E \arrow[rr,"u"] && F \arrow[ul]
\end{tikzcd}
\tag{*}
\]
is a distinguished triangle of \(\Dpx(\underline{C})\) such that \(E\) and \(F\) are perfect, \(G\)
%%%%%%%%%%  scan idx 137  |  folio 124  %%%%%%%%%%
\SGAPrintPage{124}%
is perfect. After localizing, we can suppose that \(E\) and \(F\) are objects of \(\Kpx^b(\underline{C}_0)\). After localizing further, we can suppose, by assertion (i) without the ``resp.'' alternative, that \(u\) is an arrow of \(\Kpx^b(\underline{C}_0)\); then \(G\) is perfect because it is isomorphic to the cone of \(u\), which is strictly perfect. Thus \(\Parf(\underline{C},\underline{C}_0)\) is indeed a triangulated sub-\(\underline{S}\)-category of \(\Dpx(\underline{C})\). Moreover, the preceding argument also proves the ``resp.'' assertion (ii). It remains to prove the ``resp.'' assertion (i). Let \(X\in\Ob \underline{S}\), and let \(E=(E_1\xrightarrow{a_1}E_2\xrightarrow{a_2}E_3\xrightarrow[+1]{a_3}E_1)\), \(F=(F_1\xrightarrow{b_1}F_2\xrightarrow{b_2}F_3\xrightarrow[+1]{b_3}F_1)\) be distinguished triangles of \(\Kpx^b(\underline{C}_{0X})\). Let us show that the arrow
\sgaStableTarget{sga6.e01.eq.unnamed.occ4}\[
  \underline{\Hom}_{\Tr(\Kpx^b(\underline{C}_0))}(E,F)
  \longrightarrow
  \underline{\Hom}_{\Tr(\Dpx(\underline{C}))}(E,F)
  \tag{**}
\]
is an isomorphism. By assertion (i) without the ``resp.'' alternative, we already know that \((**)\) is injective. Let \((u_i):E\to F\) be a morphism of triangles in \(\Dpx(\underline{C})\). After localizing, we can suppose, by assertion (i) without the ``resp.'' alternative, that \(u_i\) (\(1\leq i\leq3\)) is an arrow of \(\Kpx^b(\underline{C}_0)\). The arrow \(u_{i+1}a_i-b_iu_i\), for \(1\leq i\leq3\) (with the convention \(u_4=u_1\)), is zero in \(\Dpx(\underline{C})\), hence locally zero in \(\Kpx^b(\underline{C}_0)\) by assertion (i) without the ``resp.'' alternative. In other words, \((u_i)\) locally defines a morphism of triangles in \(\Kpx^b(\underline{C}_0)\). This proves that \((**)\) is surjective and completes the proof of (4.10).

\sgaStableTarget{sga6.e01.rem.4-10-1}\textbf{Remark 4.10.1.} In (4.10), we can replace \(\Parf(\underline{C},\underline{C}_0)\) by \(\Dpx^*(\underline{C})_{\underline{C}_0\text{-}\mathrm{parf}}=\Dpx^*(\underline{C})\cap\Parf(\underline{C},\underline{C}_0)\) \((*=-\text{ or }b)\); nothing in the proof needs to be changed.

\sgaStableTarget{sga6.e01.complement.4-11}\textbf{Complement 4.11.} We can describe more precisely the triangulated structure of \(\Parf(\underline{C},\underline{C}_0)\) in terms of perfect amplitude (4.7). When \(E\in\Ob\Dpx(\underline{C})\) has perfect amplitude contained in \([a,b]\), \(E[n]\) has perfect amplitude contained in \([a-n,b-n]\).

%%%%%%%%%%  scan idx 138  |  folio 125  %%%%%%%%%%
\SGAPrintPage{125}%
Let
\[
\begin{tikzcd}
& G \arrow[dl,"+1"'] & \\
E \arrow[rr] && F \arrow[ul]
\end{tikzcd}
\]
be a distinguished triangle of \(\Dpx(\underline{C})\). We can draw up a table analogous to (2.6), in which the entries indicate perfect amplitudes and each row corresponds to an implication whose conclusion is underlined.
\[
\begin{array}{ccc}
E&F&G\\[0.3em]
{}[a+1,b+1]&{}[a,b]&\underline{[a,b]}\\
{}[a,b]&\underline{[a,b]}&{}[a,b]\\
\underline{[a,b]}&{}[a,b]&{}[a-1,b-1].
\end{array}
\]
From (4.10) we immediately deduce the

\sgaStableTarget{sga6.e01.cor.4-12}\textbf{Corollary 4.12.} \uline{Suppose that} \(\underline{S}\) \uline{is a punctual site. For} \(F\in\Ob\Dpx(\underline{C})\) \uline{to have perfect amplitude (rel.}\ to \(\underline{C}_0\)\uline{) contained in} \([a,b]\)\uline{, it is necessary and sufficient that there exist a quasi-isomorphism} \(F'\to F\)\uline{, with} \(F'\) \uline{strictly perfect and such that} \(F'^i=0\) \uline{for} \(i\notin[a,b]\)\uline{. The canonical functors}
\[
  \Kpx^b(\underline{C}_0)\longrightarrow\Parf(\underline{C},\underline{C}_0)
\]
\uline{and}
\[
  \Tr(\Kpx^b(\underline{C}_0))\longrightarrow\Tr(\Parf(\underline{C},\underline{C}_0))
\]
\uline{are equivalences of categories}.

\sgaStableTarget{sga6.e01.rem.4-12-1}\textbf{Remark 4.12.1.} When \(\underline{S}\) is a punctual site, the objects of \(\underline{C}_0\) are projective. When \(\underline{C}_0\) contains all projective objects of \(\underline{C}\), one says ``projective amplitude'' instead of ``\(\underline{C}_0\)-perfect amplitude''.

% Additional notation for Expose I continuation
\providecommand{\Cpx}{\mathrm C}
\providecommand{\Kpx}{\mathrm K}
\providecommand{\Dpx}{\mathrm D}
\providecommand{\Ho}{\mathrm H}
\providecommand{\RHom}{\mathrm R\!\operatorname{Hom}}
\renewcommand{\Ob}{\operatorname{ob}}
\providecommand{\qis}{\xrightarrow{\sim}}
\providecommand{\RHom}{\mathrm R\!\Hom}
\providecommand{\Ext}{\operatorname{Ext}}
\providecommand{\Hom}{\operatorname{Hom}}
\providecommand{\Tr}{\operatorname{Tr}}

\providecommand{\Add}{\operatorname{Add}}
\providecommand{\Cart}{\underline{\operatorname{Cart}}}
\providecommand{\perfamp}{\operatorname{parf.amp}}
\providecommand{\perfdim}{\operatorname{parf.dim}}
\providecommand{\toramp}{\operatorname{tor.amp}}
\providecommand{\torf}{\operatorname{torf}}
\providecommand{\injf}{\operatorname{injf}}
\providecommand{\rg}{\operatorname{rg}}
\providecommand{\is}{\mathrm{is}}
\providecommand{\Mod}{\operatorname{Mod}}
\sgaStableTarget{sga6.e01.expose.occ3}%
